\clearpage
\chapter*{第236次报告勘误与补遗}
\markboth{第236次报告勘误与补遗}{第236次报告勘误与补遗}
\addcontentsline{toc}{chapter}{第236次报告勘误与补遗}
\sid{e236}
\src{303}

\noindent\textbf{8. TDTE VI：Picard 概形。一般性质。}

\begin{center}
[布尔巴基讨论班，第14卷，1961/62年度，第236号，23页。]
\end{center}

\noindent\textbf{第12页，第16行。} --- 把“仿射直线”改为“去掉原点的仿射
直线”。

\noindent\textbf{第18行。} --- 把“$X[t]$”改为“$X[t,t^{-1}]$”。

\noindent\textbf{第16页，命题3.1，第6行。} --- 把“然而，假设……”改为
“然而，由假设……”。

\noindent\textbf{第21页，第4节：Picard 概形的有限性定理。} --- 自本篇报告
写成以来，本节所提出的这一类有限性问题已几乎全部解决。下面指出目前在这一
方向上已知的主要事实。为简化陈述，我们默认其中出现的所有 Picard 预概形均
存在，尽管对这些陈述作显然的修改便可去掉任何存在性假设。以下用 $S$ 表示
Noether 概形，用 $X,Y$ 表示 $S$ 上的固有概形。

\medskip
\noindent(i) 设 $f:X\longrightarrow Y$ 是满射态射，则
\[
f^*: \operatorname{Pic}_{Y/S}\longrightarrow\operatorname{Pic}_{X/S}
\]
是有限型态射。

\src{304}
\noindent(ii) 假设 $Y$ 在 $S$ 上射影，并配有一个相对于 $S$ 丰沛的可逆模；
设 $X$ 是该模某一任意截面的零点概形，$f:X\longrightarrow Y$ 是典范浸入。
再假设 $Y/S$ 各纤维的不可约分支维数均为 $\geq3$。则
\[
f^*: \operatorname{Pic}_{Y/S}\longrightarrow\operatorname{Pic}_{X/S}
\]
是有限型态射。

\noindent(iii) 假设 $X$ 在 $S$ 上射影，且其所有几何纤维均为维数 $n$ 的
整概形。设 $\mathcal O_X(1)$ 是 $X$ 上的丰沛可逆模，可用来定义 Hilbert
多项式。设 $M$ 是 $\operatorname{Pic}_{X/S}$ 的一个子集；则 $M$ 拟紧，
当且仅当 $M$ 中各元素的 Hilbert 多项式
\[
a_0x^n+a_1x^{n-1}+a_2x^{n-2}+\cdots
\]
中的系数 $a_1$ 和 $a_2$ 保持有界，其中要略去单项式次数为负的系数；因此当 $n=1$ 时只出现 $a_1$，当 $n=0$ 时二者都不出现。

\noindent(iv) 对每个整数 $n\neq0$，$n$ 次幂同态在群预概形
$\operatorname{Pic}_{X/S}$ 上都是有限型态射。

注意，(i) 和 (ii) 还意味着：在这两个陈述各自的假设下，一个子集 $M$ 在
$\operatorname{Pic}_{Y/S}$ 中拟紧，当且仅当它在
$\operatorname{Pic}_{X/S}$ 中的像拟紧。由此可知，一个可逆模
$\mathcal L$ 若定义在 $Y$ 上，则它与零 $\tau$-等价，当且仅当它在 $X$ 上
的逆像具有这一性质；
换言之，$\operatorname{Pic}_{Y/S}^\tau$ 是
$\operatorname{Pic}_{X/S}^\tau$ 的逆像。特别地，要证明前一个预概形在
$S$ 上为有限型，只须对后一个证明这一点，因为
$\operatorname{Pic}_{Y/S}\longrightarrow
\operatorname{Pic}_{X/S}$ 是有限型态射。于是利用 (i)、Chow 引理和 (iii)，
得到：

\medskip
\noindent(v) $\operatorname{Pic}_{X/S}^\tau$ 在 $S$ 上为有限型。

一般地，把有限态射情形的 (i) 与 (ii) 结合起来，可把大多数有限性问题归约到
$X/S$ 的几何纤维均为维数 $\leq2$ 的正规整概形的情形；往往还可把 (i) 用于
不必有限的满射态射，再利用代数曲面的奇点消解（任意特征下由 Abhyankar
证明），将问题归约到 $X/S$ 甚至光滑、从而几何纤维为二维非奇异概形的情形。
结合 (v) 以及控制非奇异射影曲面的 Néron--Severi 群之秩的 Picard--Igusa
不等式，这例如可证明 Néron 有限性定理的如下推广：

\medskip
\noindent(vi) 设 $X/S$ 在 $S$ 上固有，除此之外任意。则 $X_i/k_i$ 是
$X/S$ 的几何纤维，其 Néron--Severi 群
\[
\operatorname{Pic}_{X_i/k_i}/\operatorname{Pic}_{X_i/k_i}^\circ
\]

\src{305}
均为有限型，而且它们的秩以及挠子群的阶均保持有界。

同样归约到非奇异曲面情形，并利用该情形下的已知定理（即 Néron 有限性定理，
以及 Néron--Severi 群上的相交形式非退化这一事实），可推出：

\medskip
\noindent(vii) 设 $X$ 是代数闭域上的固有概形。则存在有限条闭整曲线
$C_i$（$1\leq i\leq r$）位于 $X$ 中，具有如下性质：对任一子集 $M$
（取自 $\operatorname{Pic}_{X/k}$），$M$ 拟紧，当且仅当整数
\[
\deg_{C_i'}\mathcal L\qquad(\mathcal L\in M)
\]
保持有界（其中 $C_i'$ 表示 $C_i$ 的正规化）。

这里可取 $r$ 为 Néron--Severi 群的秩。一旦知道该群为有限型，(vii) 就归结为
如下事实：由曲线 $C$（位于 $X$ 中）在 Néron--Severi 群上定义的线性形式，
只在该群
的挠元素上同时为零。当 $X$ 射影且非奇异时，这一结果以及 (v) 均由
Matsusaka 得到。利用 (vii)，容易得到与零 $\tau$-等价的可逆模在 $X$ 上的
如下刻画：

\medskip
\noindent(viii) 设 $X/k$ 是域上的固有概形，$\mathcal L$ 是 $X$ 上的可逆
模。下列条件等价：

\noindent a. $\mathcal L$ 与零 $\tau$-等价；

\noindent b. 对任意凝聚模 $F$，都有
\[
\chi(F\otimes\mathcal L)=\chi(F),
\]
其中 $\chi$ 表示 Euler--Poincaré 特征；

\noindent b'. 与 (b) 相同，但取 $F=\mathcal O_Y$，其中 $Y$ 是 $X$ 中的
一维闭整子概形；

\noindent c. 对每个这样的 $Y$，以 $Y'$ 表示其正规化曲线，则
\[
\deg\mathcal L_{Y'}=0.
\]

若 $X/k$ 射影，并配有丰沛可逆模 $\mathcal O_X(1)$，则上述条件还等价于：

\noindent d. 对每个整数 $m$，$\mathcal L^{\otimes m}(1)$ 均丰沛。

\noindent e.（若 $X$ 为整概形。）对每一对整数 $m,n$，都有
\[
\chi(\mathcal L^{\otimes m}(n))=\chi(\mathcal O_X(n))
\]

\src{306}
（即在 (b) 中取 $F=\mathcal L^{\otimes m}(n)$）。

为说明最后一个条件的充分性，应注意它意味着各
$\mathcal L^{\otimes m}$ 的 Hilbert 多项式全都相等；因此根据 Mumford
判据 (iii)，各 $\mathcal L^{\otimes m}$ 留在
$\operatorname{Pic}_{X/k}$ 的一个拟紧子集中，也就是说 (a) 成立。条件
(b)、(b')、(c) 应看作通常在非奇异射影簇上定义的数值等价概念在任意固有
概形上的变体。对于非奇异射影簇，(a) 与 (c) 的等价性当然早已为人所知
（Matsusaka）。

判据 (e) 还推出如下结果：

\medskip
\noindent(ix) 设 $f:X\longrightarrow S$ 是射影平坦态射，且其几何纤维均为
整概形。则 $\operatorname{Pic}_{X/S}^\tau$ 在
$\operatorname{Pic}_{X/S}$ 中既开又闭。

下面只对关键结果 (i)、(ii)、(iii) 的证明作几点说明。（结果 (iv) 与其他结果
稍有不同，它只需对有限、满射且纯不可分的态射使用 (i) 即可证明；更确切地
说，对 Frobenius 态射使用 (i)。）证明 (i) 时，本质上使用非平坦下降的思想
（尤其参见 TDTE I，第9页）。由于这里只追求有限性结果，缺少下降数据的有效
性判据并无妨碍。另一方面，Mumford 最近证明了 (iii) 的一个稍弱形式，即涉及
Hilbert 多项式所有系数的判据。他的论证极其简单，其思路来自 Nakai 丰沛性
判据的证明（Nakai 为非奇异曲面陈述了这一判据，Mumford 将其推广到任意射影
态射）。不过在我看来，这一论证只有在对 $X/S$ 的纤维再加上一项轻微限制时
才成立，即 Serre 条件 $(S_2)$；若几何纤维正规，则该条件满足。随后在 (ii)
的证明中使用这一受限判据：(i) 确实允许归约到 $Y/S$ 平坦且几何纤维均为正规
整概形的情形；再用 Mumford 判据，很容易归约到 $X/S$ 也满足同样条件的情形。
于是由维数假设，$Y$ 和 $X$ 的几何纤维在其闭点处的深度均为 $\geq2$，因而
可以应用 SGA 1962（第XII篇报告）所给形式的“等价判据”，完成 (ii) 的证明。
一旦得到 (i) 和 (ii)，便不难在 Mumford 判据中去掉对纤维的任何正规性假设，
并证明 (iii) 中陈述的更强形式。

\src{307}
最后还应指出，(i) 和 (ii) 的证明也表明：当 $S$ 是域 $k$ 的谱时，态射
\[
\operatorname{Pic}_{Y/k}\longrightarrow\operatorname{Pic}_{X/k}
\]
是仿射的（而不只是有限型的）。

\par\smallskip
\begingroup\footnotesize
\noindent\textit{印刷本要求系数 $a_1$ 和 $a_2$ 保持有界，未说明 $n<2$ 的情形。现行读法仅明确低维时不存在的系数的约定，不宣称印刷定理为假。}\par
\endgroup

\par
